% concatenated from 221000574.tex
\documentclass[11pt]{amsart}

% Package
\usepackage{a4wide}
\usepackage{amssymb}
\usepackage{mathrsfs}
\usepackage{enumerate}
\usepackage{esint}
\usepackage{titletoc}
\usepackage[colorlinks=true, urlcolor=blue, linkcolor=blue, citecolor=black]{hyperref}
\usepackage[nameinlink]{cleveref}

% Theoremstyle plain
\theoremstyle{plain}
\newtheorem{theorem}{Theorem}[section]
\newtheorem{lemma}[theorem]{Lemma}
\newtheorem{corollary}[theorem]{Corollary}
\newtheorem{proposition}[theorem]{Proposition}

% Theoremstyle definition
\theoremstyle{definition}
\newtheorem{definition}[theorem]{Definition}
\newtheorem{remark}[theorem]{Remark}
\newtheorem{example}[theorem]{Example}

% Equation
\numberwithin{equation}{section}

% Newcommand
\DeclareMathOperator*{\esssup}{ess\,sup}
\DeclareMathOperator*{\essinf}{ess\,inf}
\DeclareMathOperator{\osc}{osc}

% MSC2020
\makeatletter
\@namedef{subjclassname@2020}{\textup{2020} Mathematics Subject Classification}
\makeatother

%%%%%%%%%%%%%%%%%%%%%%%%%%%%%%%%%%%%%%

\title[BBM theorem for fractional Sobolev spaces with variable exponents]{Bourgain, Brezis and Mironescu theorem for fractional Sobolev spaces with variable exponents}

\author{Minhyun Kim}
\address{Fakult\"at f\"ur Mathematik, Universit\"at Bielefeld, 33615 Bielefeld, Germany}
\email{minhyun.kim@uni-bielefeld.de}

\subjclass[2020]{46E35}
\keywords{fractional Sobolev space, variable exponent, limiting embedding}
\thanks{Minhyun Kim gratefully acknowledges financial support by the German Research Foundation (GRK 2235 - 282638148).}

%%%%%%%%%%%%%%%%%%%%%%%%%%%%%%%%%%%%%%

\begin{document}

\begin{abstract}
A Bourgain--Brezis--Mironescu-type theorem for fractional Sobolev spaces with variable exponents is established for sufficiently regular functions. We prove, however, that a limiting embedding theorem for these spaces fails to hold in general.
\end{abstract}

\maketitle

%%%%%%%%%%%%%%%%%%%%%%%%%%%%%%%%%%%%%%

\section{Introduction} \label{sec:introduction}

%%%%%%%%%%%%%%%%%%%%%%%%%%%%%%%%%%%%%%

In a celebrated work \cite{BBM01}, Bourgain, Brezis and Mironescu study the asymptotic behavior of the fractional Sobolev seminorms when the order of differentiability approaches one. Their results are concerned with smooth bounded domains, but the same arguments work for $W^{1, p}$-extension domain. More precisely, if $\Omega \subset \mathbb{R}^{n}$ is a $W^{1, p}$-extension domain, then an equality
\begin{equation} \label{eq:BBM}
\lim_{s \nearrow 1} s(1-s) \int_{\Omega} \int_{\Omega} \frac{|u(x)-u(y)|^{p}}{|x-y|^{n+sp}} \,\mathrm{d}y \,\mathrm{d}x = K_{n, p} \int_{\Omega} |\nabla u(x)|^{p} \,\mathrm{d}x
\end{equation}
holds for any $u \in W^{1, p}(\Omega)$, where $n \in \mathbb{N}$, $p \in (1, +\infty)$ and
\begin{equation} \label{eq:BBM-const}
K_{n, p} = \frac{1}{p} \int_{\mathbb{S}^{n-1}} |\omega_{n}|^{p} \,\mathrm{d}\mathcal{H}^{n-1}(\omega) = 2\pi^{\frac{n-1}{2}} \frac{\Gamma(\frac{p+1}{2})}{\Gamma(\frac{n+p}{2})} \frac{1}{p}.
\end{equation}
The equality \eqref{eq:BBM} suggests a nonlocal approximation of the Sobolev norm and yields a limiting embedding $W^{s, p}(\Omega) \to W^{1, p}(\Omega)$ as $s \nearrow 1$. Note that $\Omega$ is allowed to be unbounded. For the second equality in \eqref{eq:BBM-const}, see, for instance, \cite[Lemma 2.1]{IN10} or \cite[Proposition 3.9]{FG20}.

In this paper, we consider the Sobolev and fractional Sobolev spaces with variable exponents. The Sobolev spaces with variable exponents $W^{1, p(\cdot)}$ appear naturally from variational problems \cite{AM01,Alk97,CPC97,FZ99,Zhi97}, modeling of the electrorheological fluids \cite{RR96,RR01,Ruz00}, thermorheological fluids \cite{AR06} and image processing \cite{AMS08,BCE+09,CLR06,LLP10}. We refer the reader to the books \cite{CUF13,DHHR11} for a comprehensive study. The fractional Sobolev spaces with variable exponents $W^{s, p(\cdot, \cdot)}$ are introduced recently in \cite{KRV17}, and studied extensively in various contexts: \cite{ABS19,ABS20,BB19,Bah18,BR18,BRW20,BT21,CK21,DPR17,HKL21,HKS19,OK21}.

This paper aims to answer a question whether a Bourgain--Brezis--Mironescu-type result is true for the fractional Sobolev spaces with variable exponents. On the one hand, the answer is affirmative in the sense that an equality similar to \eqref{eq:BBM}, which takes the variable exponents into account, holds for functions that are sufficiently regular. However, on the other hand, the answer is negative in the sense that this equality fails to holds in general for functions in Lebesgue spaces---even in Sobolev spaces---with variable exponents. It implies that a limiting embedding for the fractional Sobolev spaces with variable exponents, corresponding to $W^{s, p} \to W^{1, p}$, is not true in general.

Let $\Omega \subset \mathbb{R}^{n}$ be any open set (not necessarily bounded). Let $p: \Omega \times \Omega \to \mathbb{R}$ be a measurable function satisfying
\begin{equation} \label{eq:p-bound}
1 < p_{-} := \essinf_{x, y \in \Omega} p(x, y) \leq \esssup_{x, y \in \Omega} p(x, y) =: p_{+} < +\infty,
\end{equation}
and define $\bar{p}(x) := p(x, x)$. Let $L^{\bar{p}(\cdot)}(\Omega)$, $W^{1, \bar{p}(\cdot)}(\Omega)$ and $W^{s, p(\cdot, \cdot)}(\Omega)$ denote the Lebesgue, Sobolev and fractional Sobolev spaces with variable exponents, respectively. See \Cref{sec:preliminaries} for the precise definitions. We begin with an extension of \eqref{eq:BBM} to an equality which takes a variable exponent $p$ into account.

\begin{theorem} \label{thm:conv}
Assume that $p$ satisfies \eqref{eq:p-bound} and
\begin{equation} \label{eq:p-log-Holder}
r^{-\osc_{B_{r}(x)} p(x, \cdot)} \leq L \quad \text{for all } x \in \Omega \text{ and } 0 < r < \min \lbrace \mathrm{dist}(x, \partial \Omega), 1 \rbrace
\end{equation}
for some constant $L \geq 1$. Then for any $u \in C^{2}_{c}(\mathbb{R}^{n})$
\begin{equation} \label{eq:equality}
\lim_{s \nearrow 1} s(1-s) \int_{\Omega} \int_{\Omega} \frac{|u(x)-u(y)|^{p(x, y)}}{|x-y|^{n+sp(x, y)}} \,\mathrm{d}y \,\mathrm{d}x = \int_{\Omega} K_{n, \bar{p}(x)} |\nabla u(x)|^{\bar{p}(x)} \,\mathrm{d}x,
\end{equation}
where
\begin{equation*}
K_{n, \bar{p}(x)} = \frac{1}{\bar{p}(x)} \int_{\mathbb{S}^{n-1}} |\omega_{n}|^{p(x)} \,\mathrm{d}\mathcal{H}^{n-1}(\omega) = 2\pi^{\frac{n-1}{2}} \frac{\Gamma(\frac{\bar{p}(x)+1}{2})}{\Gamma(\frac{n+\bar{p}(x)}{2})} \frac{1}{\bar{p}(x)}.
\end{equation*}
\end{theorem}

The condition \eqref{eq:p-log-Holder} says that for a fixed point $x \in \Omega$ the function $p(x, \cdot)$ is log-H\"older continuous at $x$. We emphasize that it does not impose any regularity on $x$-variable.

\Cref{thm:conv} is not a full extension of \eqref{eq:BBM} because \eqref{eq:equality} holds for functions that are sufficiently regular. As a natural extension of Bourgain--Brezis--Mironescu theorem, one may expect that the equality \eqref{eq:equality} holds for all functions $u \in W^{1, \bar{p}(\cdot)}(\Omega)$. However, this is not true in general even when the variable exponent $p$ is smooth, which is in sharp contrast to the case when $p$ is constant.

\begin{theorem} \label{thm:failure}
There exists a smooth variable exponent $p(\cdot, \cdot)$ such that
\begin{equation*}
W^{1, \bar{p}(\cdot)}(\Omega) \not\subset W^{s, p(\cdot, \cdot)}(\Omega)
\end{equation*}
for all $s \in (0,1)$.
\end{theorem}

The following result is an immediate consequence of \Cref{thm:failure}, which shows that a limiting embedding $W^{s, p(\cdot, \cdot)} \to W^{1, \bar{p}(\cdot)}$ as $s \nearrow 1$ is not true in general.

\begin{corollary} \label{cor:failure}
There exists a smooth variable exponent $p$ such that the equality \eqref{eq:equality} does not hold for some $u \in W^{1, \bar{p}(\cdot)}(\Omega)$.
\end{corollary}

\Cref{cor:failure} tells us that the target space $W^{1, \bar{p}(\cdot)}$ is too large for \eqref{eq:equality} to be true in general. The failure of $W^{1, \bar{p}(\cdot)}$ can be explained by the following observation: let us forget about $p(\cdot, \cdot)$ and suppose that a variable exponent $\bar{p}(\cdot)$ is given. What is the natural fractional Sobolev space as an interpolation of $L^{\bar{p}(\cdot)}$ and $W^{1, \bar{p}(\cdot)}$? Since there are too many choices of $p(\cdot, \cdot)$ satisfying \eqref{eq:p-bound} and $p(x, x) = \bar{p}(x)$, there is no way to find the natural interpolation space. This is because the spaces $L^{\bar{p}(\cdot)}$ and $W^{1, \bar{p}(\cdot)}$ do not contain any information about $p(\cdot, \cdot)$.

The following result nevertheless shows that some Sobolev functions, in a smaller space than $W^{1, \bar{p}(\cdot)}$, satisfy the equality \eqref{eq:equality}.

\begin{corollary} \label{cor:conv}
Let $\Omega$ be a $W^{1, p_{+}}$- and $W^{1, p_{-}}$-extension domain. Assume that $p$ satisfies \eqref{eq:p-bound} and \eqref{eq:p-log-Holder}. Then the equality \eqref{eq:equality} holds for all $u \in W^{1, p_{+}}(\Omega) \cap W^{1, p_{-}}(\Omega)$.
\end{corollary}

It is worth mentioning the results in \cite{CFS21,FS22}, where the authors provide nonlocal approximations to the Sobolev norms with variable exponents $p(x)$. Their approximations involve a suitable class of nonlocal nonconvex energy functionals, generalizing the constant case \cite{BN20,Ngu06,Ngu08}, and their proofs make use of the Hardy--Littlewood maximal function with variable exponent $p(x)$. Since the maximal inequality with variable exponents in a modular form fails in general \cite{Ler05}, these approximations require an additional assumption $u \in W^{1, p_{+}} \cap W^{1, p_{-}}$ as well.

Let us also mention other related results in the literature. The Bourgain--Brezis--Mironescu theorem is extended to fractional Sobolev norms with L\'evy-type measures \cite{FG20,FGKV20}, Orlicz growth \cite{ACPS20a,FBS19} and anisotropic structure \cite{CDFB22}. The asymptotic behavior of the fractional Sobolev seminorms as $s \searrow 0$ is first established in \cite{MS02}, and extended those with Orlicz growth \cite{ACPS20b} and anisotropic structure \cite{CDFB22}.

The paper is organized as follows. In \Cref{sec:preliminaries}, we recall definitions of the Lebesgue, Sobolev and fractional Sobolev spaces with variable exponents. The asymptotic behavior of the fractional Sobolev seminorms with variable exponents (in a modular form), for regular functions, is proved in \Cref{sec:limit}. \Cref{sec:failure} is devoted to the proof of \Cref{thm:failure} and \Cref{cor:failure}.

%%%%%%%%%%%%%%%%%%%%%%%%%%%%%%%%%%%%%%

\section{Preliminaries} \label{sec:preliminaries}

%%%%%%%%%%%%%%%%%%%%%%%%%%%%%%%%%%%%%%

Let us briefly recall definitions of the Lebesgue, Sobolev and fractional Sobolev spaces with variable exponents. Let $\Omega \subset \mathbb{R}^{n}$ be an open set and let $q: \Omega \to \mathbb{R}$ be a measurable function satisfying
\begin{equation*}
1 < \essinf_{x \in \Omega} q(x) \leq \esssup_{x \in \Omega} q(x) < +\infty.
\end{equation*}
The space
\begin{equation*}
L^{q(\cdot)}(\Omega) = \left\lbrace u: \Omega \to \mathbb{R} \text{ measurable}: \varrho_{q(\cdot)}(u/\lambda) < +\infty \text{ for some } \lambda > 0 \right\rbrace
\end{equation*}
equipped with the norm 
\begin{equation*}
\|u\|_{q(\cdot)} := \|u\|_{L^{q(\cdot)}(\Omega)} := \inf \left\lbrace \lambda > 0: \varrho_{q(\cdot)}(u/\lambda) \leq 1 \right\rbrace
\end{equation*}
is called the {\it variable exponent Lebesgue space}, where
\begin{equation*}
\varrho_{q(\cdot)}(u) := \varrho_{L^{q(\cdot)}(\Omega)}(u) := \int_{\Omega} |u(x)|^{q(x)}\,\mathrm{d}x
\end{equation*}
is a modular. The {\it Sobolev space with variable exponent} is defined by
\begin{equation*}
W^{1, q(\cdot)}(\Omega) = \left\lbrace u \in L^{q(\cdot)}(\Omega): D_{i}u \in L^{q(\cdot)}(\Omega) \text{ for all } i = 1, \dots, n \right\rbrace
\end{equation*}
with the norm
\begin{equation*}
\|u\|_{1, q(\cdot)} := \|u\|_{W^{1, q(\cdot)}(\Omega)} := \|u\|_{q(\cdot)} + \sum_{1 \leq i \leq n} \|D_{i} u\|_{q(\cdot)}.
\end{equation*}
It is well known that $L^{q(\cdot)}(\Omega)$ and $W^{1, q(\cdot)}(\Omega)$ are Banach spaces, see \cite{CUF13,DHHR11,FZ01,KR91} for instance.

The fractional Sobolev spaces with variable exponents are first introduced in \cite{KRV17} and studied in various contexts. Let $p: \Omega \times \Omega \to \mathbb{R}$ be a measurable function satisfying \eqref{eq:p-bound} and define $\bar{p}(x) = p(x, x)$. For $s \in (0,1)$, the {\it fractional Sobolev space with variable exponent} is defined by
\begin{equation*}
W^{s, p(\cdot, \cdot)}(\Omega) = \left\lbrace u \in L^{\bar{p}(\cdot)}(\Omega): \varrho_{s, p(\cdot, \cdot)}(u/\lambda) < + \infty \text{ for some } \lambda > 0 \right\rbrace
\end{equation*}
with the norm
\begin{equation*}
\|u\|_{s, p(\cdot, \cdot)} := \|u\|_{W^{s, p(\cdot, \cdot)}(\Omega)} := \|u\|_{p(\cdot)} + [u]_{s, p(\cdot, \cdot)},
\end{equation*}
where
\begin{equation*}
[u]_{s, p(\cdot, \cdot)} := [u]_{W^{s,p(\cdot,\cdot)}(\Omega)} := \inf \left\lbrace \lambda > 0: \varrho_{s,p(\cdot,\cdot)}(u/\lambda) \leq 1 \right\rbrace
\end{equation*}
is a seminorm and
\begin{equation*}
\varrho_{s, p(\cdot, \cdot)}(u) := \varrho_{W^{s,p(\cdot,\cdot)}(\Omega)}(u) := s(1-s) \int_{\Omega} \int_{\Omega} \frac{|u(x)-u(y)|^{p(x,y)}}{|x-y|^{n+sp(x,y)}} \,\mathrm{d}y\,\mathrm{d}x.
\end{equation*}
It is easy to check that $W^{s, p(\cdot, \cdot)}(\Omega)$ is a Banach space by following the arguments in \cite[Proposition 4.24]{DD12}. Obviously, we have $W^{1, p_{+}}(\Omega) \cap W^{1, p_{-}}(\Omega) \subset W^{s, p_{+}}(\Omega) \cap W^{s, p_{-}}(\Omega) \subset W^{s, p(\cdot, \cdot)}(\Omega)$. Let us collect some useful properties for the the fractional Sobolev spaces with variables exponents. They follow immediately from the definition.

\begin{lemma} \label{lem:modular1}
Let $u \in W^{s, p(\cdot, \cdot)}(\Omega)$, then
\begin{equation*}
\min\left\lbrace [u]_{s, p(\cdot, \cdot)}^{p_{+}}, [u]_{s, p(\cdot, \cdot)}^{p_{-}} \right\rbrace \leq \varrho_{s, p(\cdot, \cdot)}(u) \leq \max\left\lbrace [u]_{s, p(\cdot, \cdot)}^{p_{+}}, [u]_{s, p(\cdot, \cdot)}^{p_{-}} \right\rbrace.
\end{equation*}
\end{lemma}

\begin{lemma} \label{lem:modular2}
Let $u, u_{k} \in W^{s, p(\cdot, \cdot)}(\Omega)$ for $k \in \mathbb{N}$. Then the following are equivalent:
\begin{enumerate}[(i)]
\item
$[u_{k}-u]_{s, p(\cdot, \cdot)} \to 0$ as $k \to \infty$,
\item
$\varrho_{s, p(\cdot, \cdot)}(u_{k}-u) \to 0$ as $k \to \infty$.
\end{enumerate}
In particular, if one of the assertions is satisfied, then $\varrho_{s, p(\cdot, \cdot)}(u_{k}) \to \varrho_{s, p(\cdot, \cdot)}(u)$ as $k \to \infty$.
\end{lemma}

%%%%%%%%%%%%%%%%%%%%%%%%%%%%%%%%%%%%%%

\section{Limiting behavior of norms for regular functions} \label{sec:limit}

%%%%%%%%%%%%%%%%%%%%%%%%%%%%%%%%%%%%%%

In this section, we study the limiting behavior of fractional Sobolev seminorms  with variable exponents (in a modular form) for sufficiently regular functions. Let us first prove \Cref{thm:conv}, which follows from the following two lemmas: \Cref{lem:liminf} and \Cref{lem:limsup}. Throughout the paper, we denote by $C$ strictly positive constants whose exact values are not important. These constants might change from line to line.

\begin{lemma} \label{lem:liminf}
Let $\Omega \subset \mathbb{R}^{n}$ be open and assume that $p$ satisfies \eqref{eq:p-bound} and \eqref{eq:p-log-Holder}. Then
\begin{equation} \label{eq:liminf}
\int_{\Omega} K_{n, \bar{p}(x)} |\nabla u(x)|^{\bar{p}(x)} \,\mathrm{d}x \leq \liminf_{s \nearrow 1} s(1-s) \int_{\Omega} \int_{\Omega} \frac{|u(x)-u(y)|^{p(x, y)}}{|x-y|^{n+sp(x, y)}} \,\mathrm{d}y \,\mathrm{d}x
\end{equation}
for any $u \in C^{2}(\overline{\Omega})$.
\end{lemma}

\begin{proof}
Let us fix $\varepsilon > 0$, $\theta_{0} \in (0,1)$, $R > 0$ and $x \in \Omega \cap B_{R}$. Since $u \in C^{2}(\overline{\Omega})$, there exists a constant $C > 0$, depending on $R$ and $\|u\|_{C^{2}(\overline{\Omega})}$, such that
\begin{equation} \label{eq:second-difference}
|u(x+h)-u(x)-\nabla u(x) \cdot h| \leq C |h|^{2}
\end{equation}
for all $|h| < \mathrm{dist}(x, \partial \Omega)$. Since $p(x, \cdot)$ is continuous at $x$, we find $r_{1} \in (0, \mathrm{dist}(x, \partial \Omega))$, depending on $\varepsilon$, $x$, $|\nabla u(x)|$ and $p$, such that
\begin{equation} \label{eq:r1}
(1-\varepsilon) |\nabla u(x)|^{p(x, x)} \leq |\nabla u(x)|^{p(x, y)} \leq (1+\varepsilon) |\nabla u(x)|^{p(x, x)}
\end{equation}
holds for all $y \in B_{r_{1}}(x)$. Moreover, by continuity of $p(x, \cdot)$ at $x$ again, there exists $r_{2} \in (0, \mathrm{dist}(x, \partial \Omega))$, depending on $\varepsilon$, $x$ and $\theta_0$, such that
\begin{equation} \label{eq:r2}
(1-\varepsilon) \theta^{p(x, x)} \leq \theta^{p(x, y)} \leq (1+ \varepsilon) \theta^{p(x, x)}
\end{equation}
holds for all $y \in B_{r_{2}}(x)$ and $\theta \in [\theta_{0}, 1]$. We set $r_{0} = \min\lbrace r_{1}, r_{2}, 1 \rbrace$ and let $r \in (0, r_{0})$.

For $x \in \Omega \cap B_{R}$ and $h \in B_{r}$, we have from \eqref{eq:second-difference}
\begin{equation*}
|\nabla u(x) \cdot h| \leq |u(x+h)-u(x)| + C |h|^2.
\end{equation*}
Using an inequality
\begin{equation} \label{eq:alg-ineq}
(a+b)^{q} \leq (1+\varepsilon)a^{q} + C(q, \varepsilon) b^{q}, \quad a, b \geq 0, q > 1,
\end{equation}
we obtain
\begin{equation*}
|\nabla u(x) \cdot h|^{p(x, x+h)} \leq (1+\varepsilon) |u(x+h)-u(x)|^{p(x, x+h)} + C |h|^{2p(x, x+h)}
\end{equation*}
for some $C = C(R, \|u\|_{C^{2}(\overline{\Omega})}, p_{+}, p_{-}, \varepsilon) > 0$. Multiplying by $|h|^{-n-sp(x, x+h)}$ and integrating over $B_{r}$ with respect to $h$, we have
\begin{equation} \label{eq:liminf-I}
\begin{split}
I
&:= \int_{B_{r}} \frac{|\nabla u(x) \cdot h|^{p(x, x+h)}}{|h|^{n+sp(x,x+h)}} \,\mathrm{d}h \\
&\leq (1+\varepsilon) \int_{B_{r}} \frac{|u(x+h)-u(x)|^{p(x, x+h)}}{|h|^{n+sp(x, x+h)}} \,\mathrm{d}h + C \int_{B_{r}} \frac{|h|^{2p(x, x+h)}}{|h|^{n+sp(x, x+h)}} \,\mathrm{d}h
\end{split}
\end{equation}
for each $x \in \Omega \cap B_{R}$.

Let us estimate $I$ from below. By using the spherical coordinates, we write
\begin{equation*}
I = \int_{0}^{r} \int_{\mathbb{S}^{n-1}} |\nabla u(x) \cdot \omega|^{p(x, x+\rho \omega)} \rho^{-1+(1-s)p(x, x+\rho \omega)} \,\mathrm{d}\mathcal{H}^{n-1}(\omega) \,\mathrm{d}\rho.
\end{equation*}
We use the change of variables $\omega = A\xi$, where $A \in O(n)$ is a rotation satisfying $\nabla u(x) \cdot \omega = |\nabla u(x)| \xi_{n}$, so that we have
\begin{equation} \label{eq:I}
I = \int_{0}^{r} \int_{\mathbb{S}^{n-1}} |\nabla u(x)|^{p(x, x+\rho A\xi)} |\xi_{n}|^{p(x, x+\rho A \xi)} \rho^{-1+(1-s)p(x, x+\rho A\xi)} \,\mathrm{d}\mathcal{H}^{n-1}(\xi) \,\mathrm{d}\rho.
\end{equation}
Using \eqref{eq:r1} and \eqref{eq:r2}, we obtain
\begin{equation*}
I \geq (1-\varepsilon)^2 \left( \int_{0}^{r} \int_{\lbrace |\xi_{n}| \geq \theta_{0} \rbrace} |\xi_{n}|^{\bar{p}(x)} \rho^{-1+(1-s)p(x, x+\rho A\xi)} \,\mathrm{d}\mathcal{H}^{n-1}(\xi) \,\mathrm{d}\rho \right) |\nabla u(x)|^{\bar{p}(x)}.
\end{equation*}
Moreover, since the assumption \eqref{eq:p-log-Holder} yields
\begin{equation} \label{eq:log-Holder}
L^{-1} \leq \min\lbrace 1, \rho^{\osc_{B_{\rho}(x)}p(x, \cdot)} \rbrace \leq \rho^{p(x, x+\rho A\xi)-\bar{p}(x)} \leq \max\lbrace 1, \rho^{-\osc_{B_{\rho}(x)}p(x, \cdot)}\rbrace \leq L,
\end{equation}
$I$ can be further estimated as
\begin{equation*}
\begin{split}
I
&\geq \frac{(1-\varepsilon)^2}{L^{1-s}} \left( \int_{0}^{r} \int_{\lbrace |\xi_{n}| \geq \theta_{0} \rbrace} |\xi_{n}|^{\bar{p}(x)} \rho^{-1+(1-s) \bar{p}(x)} \,\mathrm{d}\mathcal{H}^{n-1}(\xi) \,\mathrm{d}\rho \right) |\nabla u(x)|^{\bar{p}(x)} \\
&\geq \frac{(1-\varepsilon)^2}{(1-s)L^{1-s}} \left( \int_{\lbrace |\xi_{n}| \geq \theta_{0} \rbrace} |\xi_{n}|^{\bar{p}(x)} \,\mathrm{d}\mathcal{H}^{n-1}(\xi) \right) \frac{r^{(1-s)\bar{p}(x)}}{\bar{p}(x)} |\nabla u(x)|^{\bar{p}(x)}.
\end{split}
\end{equation*}
Since
\begin{equation*}
\begin{split}
\int_{\lbrace |\xi_{n}| \geq \theta_{0} \rbrace} |\xi_{n}|^{\bar{p}(x)} \,\mathrm{d}\mathcal{H}^{n-1}(\xi)
&\geq \int_{\mathbb{S}^{n-1}} |\xi_{n}|^{\bar{p}(x)} \,\mathrm{d}\mathcal{H}^{n-1}(\xi) - \int_{\lbrace |\xi_{n}| < \theta_{0} \rbrace} \theta_{0}^{p_{-}} \,\mathrm{d}\mathcal{H}^{n-1}(\xi) \\
&\geq \bar{p}(x) K_{n, \bar{p}(x)} - \theta_{0} |\mathbb{S}^{n-1}|,
\end{split}
\end{equation*}
we arrive at
\begin{equation} \label{eq:liminf-I-lower}
I \geq \frac{(1-\varepsilon)^2}{(1-s)L^{1-s}} \left( \bar{p}(x) K_{n, \bar{p}(x)} - \theta_0 |\mathbb{S}^{n-1}| \right) \frac{r^{(1-s)\bar{p}(x)}}{\bar{p}(x)} |\nabla u(x)|^{\bar{p}(x)}.
\end{equation}

For the last term on the right-hand side of \eqref{eq:liminf-I}, we use the assumption \eqref{eq:p-bound} to have
\begin{equation} \label{eq:liminf-I-upper}
\int_{B_{r}} |h|^{-n+(2-s)p(x, x+h)} \,\mathrm{d}h \leq \int_{B_{r}} |h|^{-n+(2-s)p_{-}} \,\mathrm{d}h \leq |\mathbb{S}^{n-1}| \frac{r^{(2-s)p_{-}}}{(2-s)p_{-}}.
\end{equation}
Therefore, combining \eqref{eq:liminf-I}, \eqref{eq:liminf-I-lower}, \eqref{eq:liminf-I-upper} and integrating over $\Omega \cap B_{R}$ yield
\begin{equation*}
\begin{split}
&(1-\varepsilon)^{2} \frac{s}{L^{1-s}} \int_{\Omega \cap B_{R}} \left( \bar{p}(x) K_{n, \bar{p}(x)} - \theta_{0} |\mathbb{S}^{n-1}| \right) \frac{r^{(1-s)\bar{p}(x)}}{\bar{p}(x)} |\nabla u(x)|^{\bar{p}(x)} \,\mathrm{d}x \\
&\leq (1+\varepsilon) s(1-s) \int_{\Omega} \int_{\Omega} \frac{|u(x)-u(y)|^{p(x, y)}}{|x-y|^{n+sp(x, y)}} \,\mathrm{d}y \,\mathrm{d}x + C \frac{s(1-s)}{2-s} r^{(2-s)p_{-}}
\end{split}
\end{equation*}
for some $C > 0$ independent of $s$. We take $\liminf_{s \nearrow 1}$ on both sides to deduce
\begin{equation*}
\begin{split}
&(1-\varepsilon)^{2} \int_{\Omega \cap B_{R}} \left( \bar{p}(x) K_{n, \bar{p}(x)} - \theta_{0} |\mathbb{S}^{n-1}| \right) \frac{1}{\bar{p}(x)} |\nabla u(x)|^{\bar{p}(x)} \,\mathrm{d}x \\
&\leq (1+\varepsilon) \liminf_{s \nearrow 1} s(1-s) \int_{\Omega} \int_{\Omega} \frac{|u(x)-u(y)|^{p(x, y)}}{|x-y|^{n+sp(x, y)}} \,\mathrm{d}y \,\mathrm{d}x.
\end{split}
\end{equation*}
Since $\varepsilon$, $\theta_0$ and $R$ are arbitrarily chosen, we conclude \eqref{eq:liminf}.
\end{proof}

\begin{lemma} \label{lem:limsup}
Let $\Omega \subset \mathbb{R}^{n}$ be open and assume that $p$ satisfies \eqref{eq:p-bound} and \eqref{eq:p-log-Holder}. Then
\begin{equation} \label{eq:limsup}
\limsup_{s \nearrow 1} s(1-s) \int_{\Omega} \int_{\Omega} \frac{|u(x)-u(y)|^{p(x, y)}}{|x-y|^{n+sp(x, y)}} \,\mathrm{d}y \,\mathrm{d}x \leq \int_{\Omega} K_{n, \bar{p}(x)} |\nabla u(x)|^{\bar{p}(x)} \,\mathrm{d}x
\end{equation}
for any $u \in C^{2}_{c}(\mathbb{R}^n)$.
\end{lemma}

To prove \Cref{lem:limsup}, we use the dominated convergence theorem. For this purpose, we first prove the following lemma.

\begin{lemma} \label{lem:DCT}
Assume that $p$ satisfies \eqref{eq:p-bound}. Let $s_{0} \in (0,1)$ and $u \in C^{1}_{c}(\mathbb{R}^{n})$. Then, there exists a function $F \in L^{1}(\Omega)$ such that
\begin{equation} \label{eq:Fs}
F_{s}(x):= s(1-s) \int_{\Omega} \frac{|u(x)-u(y)|^{p(x, y)}}{|x-y|^{n+sp(x, y)}} \,\mathrm{d}y \leq F(x), \quad x \in \Omega,
\end{equation}
for all $s \in [s_{0} ,1)$.
\end{lemma}

\begin{proof}
Let us write $F_{s}(x) = I_{1} + I_{2}$, where
\begin{equation*}
\begin{split}
I_{1} &= s(1-s) \int_{\Omega \cap B_{1}(x)} \frac{|u(x)-u(y)|^{p(x, y)}}{|x-y|^{n+sp(x, y)}} \,\mathrm{d}y \quad\text{and}\\
I_{2} &= s(1-s) \int_{\Omega \setminus B_{1}(x)} \frac{|u(x)-u(y)|^{p(x, y)}}{|x-y|^{n+sp(x, y)}} \,\mathrm{d}y.
\end{split}
\end{equation*}
Using
\begin{equation*}
|u(y)-u(x)| \leq \int_{0}^{1} |\nabla u(x+t(y-x)) \cdot (y-x)| \,\mathrm{d}t \leq \|\nabla u\|_{L^{\infty}(\mathbb{R}^{n})} |y-x|
\end{equation*}
and \eqref{eq:p-bound}, we have
\begin{equation} \label{eq:DCT-I1}
\begin{split}
I_{1}
&\leq (1-s) \int_{\Omega \cap B_{1}(x)} \|\nabla u\|_{L^{\infty}(\mathbb{R}^{n})}^{p(x, y)} |x-y|^{-n+(1-s)p(x, y)} \,\mathrm{d}y \\
&\leq (1-s) \left( \|\nabla u\|_{L^{\infty}(\mathbb{R}^{n})}^{p_{+}} +  \|\nabla u\|_{L^{\infty}(\mathbb{R}^{n})}^{p_{-}} \right) \int_{B_{1}} |h|^{-n+(1-s)p_{-}} \,\mathrm{d}h \\
&= \frac{|\mathbb{S}^{n-1}|}{p_{-}} \left( \|\nabla u\|_{L^{\infty}(\mathbb{R}^{n})}^{p_{+}} +  \|\nabla u\|_{L^{\infty}(\mathbb{R}^{n})}^{p_{-}} \right).
\end{split}
\end{equation}
Moreover, using \eqref{eq:p-bound} again, we obtain
\begin{equation} \label{eq:DCT-I2}
\begin{split}
I_{2}
&\leq 2^{p_+} s \int_{\Omega \setminus B_{1}(x)} \|u\|_{L^{\infty}(\mathbb{R}^{n})}^{p(x, x+h)} |x-y|^{-n-sp(x, y)} \,\mathrm{d}y \\
&\leq 2^{p_+} s \left( \|u\|_{L^{\infty}(\mathbb{R}^{n})}^{p_{+}} + \|u\|_{L^{\infty}(\mathbb{R}^{n})}^{p_{-}} \right) \int_{\mathbb{R}^{n} \setminus B_{1}} |h|^{-n-sp_{-}} \,\mathrm{d}h \\
&= 2^{p_{+}} \frac{|\mathbb{S}^{n-1}|}{p_{-}} \left( \|u\|_{L^{\infty}(\mathbb{R}^{n})}^{p_{+}} + \|u\|_{L^{\infty}(\mathbb{R}^{n})}^{p_{-}} \right).
\end{split}
\end{equation}
Thus, it follows from \eqref{eq:DCT-I1} and \eqref{eq:DCT-I2}
\begin{equation} \label{eq:Fs-bound}
F_{s} = I_{1} + I_{2} \leq C \quad\text{in } \Omega
\end{equation}
for all $s \in (0,1)$, where $C$ is a constant depending on $n$, $p_{+}$, $p_{-}$ and $\|u\|_{C^1(\mathbb{R}^{n})}$. This finishes the proof when $\Omega$ is bounded.

Let us next consider the case when $\Omega$ is bounded. We prove that $F_{s}$ is bounded by an integrable function for all $s \in [s_{0},1)$. To this end, let us consider a large ball $B_{R/2}$ containing $\mathrm{supp} \,u$. We may assume that $R > 2$. Then, for $x \in \Omega \setminus B_{R}$, we have
\begin{equation*}
F_{s}(x) = s(1-s) \int_{\Omega \cap B_{R/2}} \frac{|u(y)|^{p(x, y)}}{|x-y|^{n+sp(x, y)}} \,\mathrm{d}y.
\end{equation*}
Since $|x-y| \geq |x|-R/2 \geq |x|/2 \geq 1$, using \eqref{eq:p-bound} we obtain
\begin{equation} \label{eq:Fs-decay}
F_{s}(x) \leq |B_{R/2}| \left( \|u\|_{L^{\infty}(\mathbb{R}^{n})}^{p_{+}} + \|u\|_{L^{\infty}(\mathbb{R}^{n})}^{p_{-}} \right) \left( \frac{2}{|x|} \right)^{n+s_{0} p_{-}} \quad\text{for } x \in \Omega \setminus B_{R}.
\end{equation}
Combining \eqref{eq:Fs-bound} and \eqref{eq:Fs-decay}, we conclude that
\begin{equation*}
F_{s}(x) \leq C \min \left\lbrace \frac{1}{R^{n+s_{0}p_{-}}}, \frac{1}{|x|^{n+s_{0}p_{-}}} \right\rbrace =: F(x) \in L^{1}(\Omega),
\end{equation*}
where $C$ is a constant depending on $n$, $p_{+}$, $p_{-}$, $s_{0}$ and $u$.
\end{proof}

In the next lemma, we find a pointwise limit of $F_{s}$ as $s \nearrow 1$.

\begin{lemma} \label{lem:limit-Fs}
Assume that $p$ satisfies \eqref{eq:p-bound} and \eqref{eq:p-log-Holder}. Let $u \in C^{2}_{c}(\mathbb{R}^{n})$ and let $F_{s}$ be given as in \eqref{eq:Fs}. Then
\begin{equation*}
\lim_{s \nearrow 1} F_{s}(x) = K_{n, \bar{p}(x)} |\nabla u(x)|^{\bar{p}(x)}
\end{equation*}
for all $x \in \Omega$.
\end{lemma}

\begin{proof}
Let $u \in C^{2}_{c}(\mathbb{R}^{n})$, then there exists a constant $C > 0$, depending on $\|u\|_{C^{2}(\mathbb{R}^{n})}$, such that
\begin{equation} \label{eq:gradient}
|u(y)-u(x)| \leq |\nabla u(x) \cdot (y-x)| + C|y-x|^2
\end{equation}
holds for all $x, y \in \mathbb{R}^{n}$. Let $\varepsilon > 0$ and $\theta_{0} \in (0,1)$. Similarly as in the proof of \Cref{lem:liminf}, for a fixed point $x \in \Omega$ we find $r_{0} \in (0,1)$ such that \eqref{eq:r1} and \eqref{eq:r2} hold for all $y \in B_{r_{0}}(x)$. Then, using \eqref{eq:gradient} and \eqref{eq:alg-ineq} we have
\begin{equation} \label{eq:limsup-I12}
\begin{split}
\int_{\Omega} \frac{|u(x)-u(y)|^{p(x, y)}}{|x-y|^{n+sp(x, y)}} \,\mathrm{d}y
&\leq (1+\varepsilon) \int_{B_{r}(x)} \frac{|\nabla u(x) \cdot (y-x)|^{p(x, y)}}{|x-y|^{n+sp(x, y)}} \,\mathrm{d}y \\
&\quad + C \int_{B_{r}(x)} \frac{|x-y|^{2p(x, y)}}{|x-y|^{n+sp(x, y)}} \,\mathrm{d}y \\
&\quad + \int_{\Omega \setminus B_{r}(x)} \frac{|u(x)-u(y)|^{p(x, y)}}{|x-y|^{n+sp(x, y)}} \,\mathrm{d}y \\
&=: I_{1} + I_{2} + I_{3}
\end{split}
\end{equation}
for $r \in (0, r_{0})$.

Let us estimate $I_{1}$, $I_{2}$ and $I_{3}$. We first recall from \eqref{eq:I} and \eqref{eq:liminf-I-upper} that
\begin{equation*}
I_{1} = (1+\varepsilon) \int_{0}^{r} \int_{\mathbb{S}^{n-1}} |\nabla u(x)|^{p(x, x+\rho A\xi)} |\xi_{n}|^{p(x, x+\rho A \xi)} \rho^{-1+(1-s)p(x, x+\rho A\xi)} \,\mathrm{d}\mathcal{H}^{n-1}(\xi) \,\mathrm{d}\rho
\end{equation*}
and
\begin{equation} \label{eq:limsup-I2}
I_{2} \leq C \frac{r^{(2-s)p_-}}{2-s}
\end{equation}
hold. For $I_{1}$, we use \eqref{eq:r1} and \eqref{eq:log-Holder} to have
\begin{equation*}
I_{1} \leq (1+\varepsilon)^{2} L^{1-s} \left( \int_{0}^{r} \int_{\mathbb{S}^{n-1}} |\xi_{n}|^{p(x, x+\rho A\xi)} \,\mathrm{d}\mathcal{H}^{n-1}(\xi) \,\rho^{-1+(1-s)\bar{p}(x)} \,\mathrm{d}\rho \right) |\nabla u(x)|^{\bar{p}(x)}.
\end{equation*}
Moreover, by using \eqref{eq:r2} we obtain
\begin{equation*}
\begin{split}
\int_{\mathbb{S}^{n-1}} |\xi_{n}|^{p(x, x+\rho A\xi)} \,\mathrm{d}\mathcal{H}^{n-1}(\xi)
&\leq \int_{\lbrace |\xi_{n}| \geq \theta_{0} \rbrace} |\xi_n|^{p(x, x+\rho A\xi)} \,\mathrm{d}\mathcal{H}^{n-1}(\xi) + \theta_{0} |\mathbb{S}^{n-1}| \\
&\leq (1+\varepsilon) \bar{p}(x) K_{n, \bar{p}(x)} + \theta_{0} |\mathbb{S}^{n-1}|.
\end{split}
\end{equation*}
Thus, we deduce
\begin{equation} \label{eq:limsup-I1}
I_{1} \leq (1+\varepsilon)^{2} \frac{L^{1-s}}{1-s} \left( (1+\varepsilon) \bar{p}(x) K_{n, \bar{p}(x)} + \theta_{0} |\mathbb{S}^{n-1}| \right) \frac{r^{(1-s)\bar{p}(x)}}{\bar{p}(x)} |\nabla u(x)|^{\bar{p}(x)}.
\end{equation}
For $I_{3}$, it follows from $|x-y|^{-sp(x, y)} \leq r^{sp_{-} - sp(x, y)} |x-y|^{-sp_{-}} \leq r^{s(p_{-} - p_{+})} |x-y|^{-sp_{-}}$
\begin{equation} \label{eq:limsup-I3}
I_{3} \leq \frac{2^{p_{+}} |\mathbb{S}^{n-1}|}{sp_{-}} \left( \|u\|_{L^{\infty}(\mathbb{R}^{n})}^{p_{+}} + \|u\|_{L^{\infty}(\mathbb{R}^{n})}^{p_{-}} \right) r^{-sp_{+}}.
\end{equation}
Combining \eqref{eq:limsup-I12}--\eqref{eq:limsup-I3}, we arrive at
\begin{equation*}
\begin{split}
F_{s}(x)
&= s(1-s) (I_{1} + I_{2} + I_{3}) \\
&\leq (1+\varepsilon)^{2} sL^{1-s} \left( (1+\varepsilon) \bar{p}(x) K_{n, \bar{p}(x)} + \theta_{0} |\mathbb{S}^{n-1}| \right) \frac{r^{(1-s)\bar{p}(x)}}{\bar{p}(x)} |\nabla u(x)|^{\bar{p}(x)} \\
&\quad + C \frac{s(1-s)}{2-s} r^{(2-s)p_{-}} + \frac{2^{p_{+}} |\mathbb{S}^{n-1}|}{p_{-}} (1-s) \left( \|u\|_{L^{\infty}(\mathbb{R}^{n})}^{p_{+}} + \|u\|_{L^{\infty}(\mathbb{R}^{n})}^{p_{-}} \right) r^{-sp_{+}},
\end{split}
\end{equation*}
and hence
\begin{equation*}
\limsup_{s \nearrow 1} F_{s}(x) \leq (1+\varepsilon)^{2} \left( (1+\varepsilon) \bar{p}(x) K_{n, \bar{p}(x)} + \theta_{0} |\mathbb{S}^{n-1}| \right) \frac{1}{\bar{p}(x)} |\nabla u(x)|^{\bar{p}(x)}.
\end{equation*}
Since $\varepsilon$ and $\theta_{0}$ are arbitrarily chosen, we deduce $\limsup_{s \nearrow 1} F_{s}(x) \leq K_{n, \bar{p}(x)} |\nabla u(x)|^{\bar{p}(x)}$. Recall from \eqref{eq:liminf-I-lower}
\begin{equation*}
F_{s}(x) \geq (1+\varepsilon)(1-\varepsilon)^{2} \frac{s}{L^{1-s}} \left( \bar{p}(x) K_{n, \bar{p}(x)} - \theta_{0} |\mathbb{S}^{n-1}| \right) \frac{r^{(1-s)\bar{p}(x)}}{\bar{p}(x)} |\nabla u(x)|^{\bar{p}(x)},
\end{equation*}
which results in $\liminf_{s \nearrow 1} F_{s}(x) \geq K_{n, \bar{p}(x)} |\nabla u(x)|^{\bar{p}(x)}$, finishing the proof.
\end{proof}

\Cref{lem:limsup} follows from \Cref{lem:DCT} and \Cref{lem:limit-Fs}. Moreover, \Cref{thm:conv} is an immediate consequence of \Cref{lem:liminf}, \Cref{lem:limsup} and the dominated convergence theorem. Let us finish this section by providing the proof of \Cref{cor:conv}.

\begin{proof} [Proof of \Cref{cor:conv}]
Since $\Omega$ is a $W^{1, p_{+}}$- and $W^{1, p_{-}}$-extension domain, we may assume that $u \in W^{1, p_{+}}(\mathbb{R}^{n}) \cap W^{1, p_{-}}(\mathbb{R}^{n})$. There exists a sequence $\lbrace u_{k} \rbrace \subset C^{2}_{c}(\mathbb{R}^{n})$ such that $u_{k} \to u$ in $W^{1, p_{+}}(\mathbb{R}^{n})$ and $W^{1, p_{-}}(\mathbb{R}^{n})$ as $k \to \infty$. By the triangle inequality, we have
\begin{equation*}
\begin{split}
\left| \varrho_{W^{s, p(\cdot, \cdot)}(\Omega)}(u) - \tilde{\varrho}_{L^{\bar{p}(\cdot)}(\Omega)}(|\nabla u|) \right|
&\leq \left| \varrho_{W^{s, p(\cdot, \cdot)}(\Omega)}(u) - \varrho_{W^{s, p(\cdot, \cdot)}(\Omega)}(u_{k}) \right| \\
&\quad + \left| \varrho_{W^{s, p(\cdot, \cdot)}(\Omega)}(u_{k}) - \tilde{\varrho}_{L^{\bar{p}(\cdot)}(\Omega)}(|\nabla u_{k}|) \right| \\
&\quad + \left| \tilde{\varrho}_{L^{\bar{p}(\cdot)}(\Omega)}(|\nabla u_{k}|) - \tilde{\varrho}_{L^{\bar{p}(\cdot)}(\Omega)}(|\nabla u|) \right| \\
&=: I_{1} + I_{2} + I_{3},
\end{split}
\end{equation*}
where
\begin{equation*}
\tilde{\varrho}_{L^{\bar{p}(\cdot)}(\Omega)}(v) = \int_{\Omega} K_{n, \bar{p}}(x) |v(x)|^{\bar{p}(x)} \,\mathrm{d}x.
\end{equation*}
Since $W^{1, p_{+}}(\mathbb{R}^{n}) \cap W^{1, p_{-}}(\mathbb{R}^{n}) \subset W^{s, p(\cdot, \cdot)}(\mathbb{R}^{n}) \subset W^{s, p(\cdot, \cdot)}(\Omega)$, $u_{k} \to u$ in $W^{s, p(\cdot, \cdot)}(\Omega)$ as $k \to \infty$. Thus, by \Cref{lem:modular2}, we have $I_{1} \to 0$ as $k \to \infty$. Obviously, we have $I_{3} \to 0$ as $k \to \infty$. Therefore, for $\varepsilon > 0$, we obtain
\begin{equation*}
\left| \varrho_{W^{s, p(\cdot, \cdot)}(\Omega)}(u) - \tilde{\varrho}_{L^{\bar{p}(\cdot)}(\Omega)}(|\nabla u|) \right| \leq \varepsilon + I_{2}
\end{equation*}
for sufficiently large $k$. By \Cref{thm:conv} we arrive at
\begin{equation*}
\limsup_{s \nearrow 1} \left| \varrho_{W^{s, p(\cdot, \cdot)}(\Omega)}(u) - \tilde{\varrho}_{L^{\bar{p}(\cdot)}(\Omega)}(|\nabla u|) \right| \leq \varepsilon,
\end{equation*}
which conclude the theorem.
\end{proof}

%%%%%%%%%%%%%%%%%%%%%%%%%%%%%%%%%%%%%%

\section{Failure of the limiting embedding} \label{sec:failure}

%%%%%%%%%%%%%%%%%%%%%%%%%%%%%%%%%%%%%%

Having \Cref{thm:conv} at hand, one may expect the limiting embedding $W^{s, p(\cdot, \cdot)} \to W^{1, \bar{p}(\cdot)}$ as a natural extension of the limiting embedding $W^{s, p} \to W^{1, p}$. However, in this section, we provide an example of a variable exponent $p(\cdot, \cdot)$ such that the limiting embedding $W^{s, p(\cdot, \cdot)} \to W^{1, \bar{p}(\cdot)}$ fails. As explained in the introduction, the idea is to construct an unnatural fractional Sobolev space as one of the interpolations of the given Lebesgue and Sobolev spaces. Such a strange space can be constructed even from $L^{p}$ and $W^{1, p}$ with a constant $p$. In the following proof, we find a variable exponent $p(\cdot, \cdot)$ such that $\bar{p}(x) = p(x, x) \equiv p$ is a constant and that $W^{1, p} \not\subset W^{s, p(\cdot, \cdot)}$ for all $s \in (0,1)$.

\begin{proof} [Proof of \Cref{thm:failure}]
We may assume that $\Omega$ contains $\overline{B}_{2}$. The idea of proof is to construct a variable exponent $p$ so that $W^{s, p(\cdot, \cdot)}(\Omega)$ requires more integrability than $W^{1, \bar{p}(\cdot)}(\Omega)$. To this end, let $p: \Omega \times \Omega \to \mathbb{R}$ be a smooth function satisfying \eqref{eq:p-bound}. Slightly abusing notation, we assume that $p(x, y) = p(|x-y|)$ for all $x, y \in \Omega$. Then, $\bar{p} = \bar{p}(x) = p(x, x) = p(0)$ is a constant. For $q \in (1, n)$, we assume
\begin{equation} \label{eq:p}
1 < \bar{p} < n/q \quad\text{and}\quad p(r) \geq n/(q-1) \quad\text{for } r \geq 1.
\end{equation}
The assumption \eqref{eq:p} shows how $W^{s, p(\cdot, \cdot)}(\Omega)$ requires more integrability than $W^{1, \bar{p}}(\Omega)$.

We consider a smooth decreasing function $\phi: (0, +\infty) \to [0, +\infty)$ satisfying $\phi(t) = |t|^{1-q}$ for $t \in (0, 1]$, $0 \leq \phi(t) \leq 1$ for $t \in [1, 2]$ and $\phi(t) = 0$ for $t \in [2, +\infty)$. We claim that a function $u: \Omega \to [0, +\infty]$ defined by $u(x) = \phi(|x|)$ belongs to $W^{1, \bar{p}}(\Omega) \setminus W^{s, p(\cdot, \cdot)}(\Omega)$ for all $s \in (0,1)$. Indeed, a straightforward computation shows
\begin{equation*}
\int_{\Omega} |u(x)|^{\bar{p}} \,\mathrm{d}x \leq \int_{B_{1}} |x|^{-(q-1)\bar{p}} \,\mathrm{d}x + |B_{2} \setminus B_{1}| < +\infty
\end{equation*}
and
\begin{equation*}
\int_{\Omega} |\nabla u(x)|^{\bar{p}} \,\mathrm{d}x \leq (q-1)^{\bar{p}} \int_{B_{1}} |x|^{-q\bar{p}} \,\mathrm{d}x + C |B_{2} \setminus B_{1}| < + \infty.
\end{equation*}
Thus, $u \in W^{1, \bar{p}}(\Omega)$.

It remains to prove $u \not\in W^{s, p(\cdot, \cdot)}(\Omega)$ for all $s \in (0,1)$. Since $u=0$ outside $B_{2}$, we have
\begin{equation*}
\varrho_{s, p(\cdot, \cdot)}(u) \geq s(1-s) \int_{\Omega \setminus B_{2}} \int_{B_{1}} \frac{|y|^{-(q-1)p(|x-y|)}}{|x-y|^{n+sp(|x-y|)}} \,\mathrm{d}y \,\mathrm{d}x.
\end{equation*}
If $x \in \Omega \setminus B_{2}$ and $y \in B_{1}$, then $1 \leq |x-y| \leq |x|+1 \leq \frac{3}{2} |x|$. Thus, by using \eqref{eq:p-bound} and \eqref{eq:p}, we obtain
\begin{equation*}
\varrho_{s, p(\cdot, \cdot)}(u) \geq s(1-s) \int_{\Omega \setminus B_{2}} \int_{B_{1}} \frac{|y|^{-n}}{(\frac{3}{2}|x|)^{n+sp_+}} \,\mathrm{d}y \,\mathrm{d}x = +\infty.
\end{equation*}
Therefore, it follows from \Cref{lem:modular1} $u \notin W^{s, p(\cdot, \cdot)}(\Omega)$ for all $s \in (0,1)$.
\end{proof}

%%%%%%%%%%%%%%%%%%%%%%%%%%%%%%%%%%%%%%

\begin{thebibliography}{10}

\bibitem{AMS08}
R.~Aboulaich, D.~Meskine, and A.~Souissi.
\newblock New diffusion models in image processing.
\newblock {\em Comput. Math. Appl.}, 56(4):874--882, 2008.

\bibitem{AM01}
E.~Acerbi and G.~Mingione.
\newblock Regularity results for a class of functionals with non-standard
  growth.
\newblock {\em Arch. Ration. Mech. Anal.}, 156(2):121--140, 2001.

\bibitem{ACPS20a}
A.~Alberico, A.~Cianchi, L.~Pick, and L.~Slav\'{\i}kov\'{a}.
\newblock On the limit as {$s \to 1^-$} of possibly non-separable fractional
  {O}rlicz-{S}obolev spaces.
\newblock {\em Atti Accad. Naz. Lincei Rend. Lincei Mat. Appl.},
  31(4):879--899, 2020.

\bibitem{ACPS20b}
A.~Alberico, A.~Cianchi, L.~Pick, and L.~Slav\'{\i}kov\'{a}.
\newblock On the limit as {$s\to 0^+$} of fractional {O}rlicz-{S}obolev spaces.
\newblock {\em J. Fourier Anal. Appl.}, 26(6):Paper No. 80, 19, 2020.

\bibitem{Alk97}
Y.~A. Alkhutov.
\newblock The {H}arnack inequality and the {H}\"{o}lder property of solutions
  of nonlinear elliptic equations with a nonstandard growth condition.
\newblock {\em Differ. Uravn.}, 33(12):1651--1660, 1726, 1997.

\bibitem{AR06}
S.~N. Antontsev and J.~F. Rodrigues.
\newblock On stationary thermo-rheological viscous flows.
\newblock {\em Ann. Univ. Ferrara Sez. VII Sci. Mat.}, 52(1):19--36, 2006.

\bibitem{ABS19}
E.~Azroul, A.~Benkirane, and M.~Shimi.
\newblock Eigenvalue problems involving the fractional {$p(x)$}-{L}aplacian
  operator.
\newblock {\em Adv. Oper. Theory}, 4(2):539--555, 2019.

\bibitem{ABS20}
E.~Azroul, A.~Benkirane, and M.~Shimi.
\newblock General fractional {S}obolev space with variable exponent and
  applications to nonlocal problems.
\newblock {\em Adv. Oper. Theory}, 5(4):1512--1540, 2020.

\bibitem{BB19}
A.~Baalal and M.~Berghout.
\newblock Density properties for fractional {S}obolev spaces with variable
  exponents.
\newblock {\em Ann. Funct. Anal.}, 10(3):308--324, 2019.

\bibitem{Bah18}
A.~Bahrouni.
\newblock Comparison and sub-supersolution principles for the fractional
  {$p(x)$}-{L}aplacian.
\newblock {\em J. Math. Anal. Appl.}, 458(2):1363--1372, 2018.

\bibitem{BR18}
A.~Bahrouni and V.~D. R\u{a}dulescu.
\newblock On a new fractional {S}obolev space and applications to nonlocal
  variational problems with variable exponent.
\newblock {\em Discrete Contin. Dyn. Syst. Ser. S}, 11(3):379--389, 2018.

\bibitem{BRW20}
A.~Bahrouni, V.~D. R\u{a}dulescu, and P.~Winkert.
\newblock Robin fractional problems with symmetric variable growth.
\newblock {\em J. Math. Phys.}, 61(10):101503, 14, 2020.

\bibitem{BT21}
R.~Biswas and S.~Tiwari.
\newblock Variable order nonlocal {C}hoquard problem with variable exponents.
\newblock {\em Complex Var. Elliptic Equ.}, 66(5):853--875, 2021.

\bibitem{BCE+09}
E.~M. Bollt, R.~Chartrand, S.~Esedo\={g}lu, P.~Schultz, and K.~R. Vixie.
\newblock Graduated adaptive image denoising: local compromise between total
  variation and isotropic diffusion.
\newblock {\em Adv. Comput. Math.}, 31(1-3):61--85, 2009.

\bibitem{BBM01}
J.~Bourgain, H.~Brezis, and P.~Mironescu.
\newblock Another look at {S}obolev spaces.
\newblock In {\em Optimal control and partial differential equations}, pages
  439--455. IOS, Amsterdam, 2001.

\bibitem{BN20}
H.~Brezis and H.-M. Nguyen.
\newblock Non-local, non-convex functionals converging to {S}obolev norms.
\newblock {\em Nonlinear Anal.}, 191:111626, 9, 2020.

\bibitem{CK21}
J.~Chaker and M.~Kim.
\newblock Local regularity for nonlocal equations with variable exponents.
\newblock {\em arXiv preprint arXiv:2107.06043}, 2021.

\bibitem{CLR06}
Y.~Chen, S.~Levine, and M.~Rao.
\newblock Variable exponent, linear growth functionals in image restoration.
\newblock {\em SIAM J. Appl. Math.}, 66(4):1383--1406, 2006.

\bibitem{CPC97}
V.~Chiad\`o~Piat and A.~Coscia.
\newblock H\"{o}lder continuity of minimizers of functionals with variable
  growth exponent.
\newblock {\em Manuscripta Math.}, 93(3):283--299, 1997.

\bibitem{CFS21}
I.~Cinelli, G.~Ferrari, and M.~Squassina.
\newblock Nonlocal approximations to anisotropic {S}obolev norms.
\newblock {\em arXiv preprint arXiv:2110.13662}, 2021.

\bibitem{CUF13}
D.~V. Cruz-Uribe and A.~Fiorenza.
\newblock {\em Variable {L}ebesgue spaces}.
\newblock Applied and Numerical Harmonic Analysis. Birkh\"{a}user/Springer,
  Heidelberg, 2013.
\newblock Foundations and harmonic analysis.

\bibitem{DPR17}
L.~M. Del~Pezzo and J.~D. Rossi.
\newblock Traces for fractional {S}obolev spaces with variable exponents.
\newblock {\em Adv. Oper. Theory}, 2(4):435--446, 2017.

\bibitem{DD12}
F.~Demengel and G.~Demengel.
\newblock {\em Functional spaces for the theory of elliptic partial
  differential equations}.
\newblock Universitext. Springer, London; EDP Sciences, Les Ulis, 2012.
\newblock Translated from the 2007 French original by Reinie Ern\'{e}.

\bibitem{DHHR11}
L.~Diening, P.~Harjulehto, P.~H\"{a}st\"{o}, and M.~R{\r u}\v{z}i\v{c}ka.
\newblock {\em Lebesgue and {S}obolev spaces with variable exponents}, volume
  2017 of {\em Lecture Notes in Mathematics}.
\newblock Springer, Heidelberg, 2011.

\bibitem{CDFB22}
I.~Ceresa~Dussel and J.~Fern\'{a}ndez~Bonder.
\newblock A {B}ourgain-{B}rezis-{M}ironescu formula for anisotropic fractional
  {S}obolev spaces and applications to anisotropic fractional differential
  equations.
\newblock {\em arXiv preprint arXiv:2206.11873}, 2022.

\bibitem{FZ99}
X.~Fan and D.~Zhao.
\newblock A class of {D}e {G}iorgi type and {H}\"{o}lder continuity.
\newblock {\em Nonlinear Anal.}, 36(3, Ser. A: Theory Methods):295--318, 1999.

\bibitem{FZ01}
X.~Fan and D.~Zhao.
\newblock On the spaces {$L^{p(x)}(\Omega)$} and {$W^{m,p(x)}(\Omega)$}.
\newblock {\em J. Math. Anal. Appl.}, 263(2):424--446, 2001.

\bibitem{FBS19}
J.~Fern\'{a}ndez~Bonder and A.~M. Salort.
\newblock Fractional order {O}rlicz-{S}obolev spaces.
\newblock {\em J. Funct. Anal.}, 277(2):333--367, 2019.

\bibitem{FS22}
G.~Ferrari and M.~Squassina.
\newblock Nonlocal characterizations of variable exponent {S}obolev spaces.
\newblock {\em Asymptot. Anal.}, 127(1-2):121--142, 2022.

\bibitem{FG20}
G.~F. Foghem~Gounoue.
\newblock A remake on the {B}ourgain-{B}rezis-{M}ironescu characterization of {S}obolev
  spaces.
\newblock {\em arXiv preprint arXiv:2008.07631}, 2020.

\bibitem{FGKV20}
G.~F. Foghem~Gounoue, M.~Kassmann, and P.~Voigt.
\newblock Mosco convergence of nonlocal to local quadratic forms.
\newblock {\em Nonlinear Anal.}, 193:111504, 22, 2020.

\bibitem{HKL21}
K.~Ho, Y.-H. Kim, and J.~Lee.
\newblock Schr\"{o}dinger {$p(\cdot)$}-{L}aplace equations in {$\Bbb R^N$}
  involving indefinite weights and critical growth.
\newblock {\em J. Math. Phys.}, 62(11):Paper No. 111506, 28, 2021.

\bibitem{HKS19}
K.~Ho, Y.-H. Kim, and I.~Sim.
\newblock Existence results for {S}chr\"{o}dinger {$p(\cdot)$}-{L}aplace
  equations involving critical growth in {$\Bbb R^N$}.
\newblock {\em Nonlinear Anal.}, 182:20--44, 2019.

\bibitem{IN10}
H.~Ishii and G.~Nakamura.
\newblock A class of integral equations and approximation of {$p$}-{L}aplace
  equations.
\newblock {\em Calc. Var. Partial Differential Equations}, 37(3-4):485--522,
  2010.

\bibitem{KRV17}
U.~Kaufmann, J.~D. Rossi, and R.~Vidal.
\newblock Fractional {S}obolev spaces with variable exponents and fractional
  {$p(x)$}-{L}aplacians.
\newblock {\em Electron. J. Qual. Theory Differ. Equ.}, Paper No. 76, 10 pp,
  2017.

\bibitem{KR91}
O.~Kov\'{a}\v{c}ik and J.~R\'{a}kosn\'{\i}k.
\newblock On spaces {$L^{p(x)}$} and {$W^{k,p(x)}$}.
\newblock {\em Czechoslovak Math. J.}, 41(116)(4):592--618, 1991.

\bibitem{Ler05}
A.~K. Lerner.
\newblock Some remarks on the {H}ardy-{L}ittlewood maximal function on variable
  {$L^p$} spaces.
\newblock {\em Math. Z.}, 251(3):509--521, 2005.

\bibitem{LLP10}
F.~Li, Z.~Li, and L.~Pi.
\newblock Variable exponent functionals in image restoration.
\newblock {\em Appl. Math. Comput.}, 216(3):870--882, 2010.

\bibitem{MS02}
V.~Mazya and T.~Shaposhnikova.
\newblock On the {B}ourgain, {B}rezis, and {M}ironescu theorem concerning
  limiting embeddings of fractional {S}obolev spaces.
\newblock {\em J. Funct. Anal.}, 195(2):230--238, 2002.

\bibitem{Ngu06}
H.-M. Nguyen.
\newblock Some new characterizations of {S}obolev spaces.
\newblock {\em J. Funct. Anal.}, 237(2):689--720, 2006.

\bibitem{Ngu08}
H.-M. Nguyen.
\newblock Further characterizations of {S}obolev spaces.
\newblock {\em J. Eur. Math. Soc. (JEMS)}, 10(1):191--229, 2008.

\bibitem{OK21}
J.~Ok.
\newblock Local {H}\"{o}lder regularity for nonlocal equations with
  variable powers.
\newblock {\em arXiv preprint arXiv:2107.06611}, 2021.

\bibitem{RR96}
K.~R. Rajagopal and M.~R{\r u}\v{z}i\v{c}ka.
\newblock On the modeling of electrorheological materials.
\newblock {\em Mech. Research Comm.}, 23(4):401--407, 1996.

\bibitem{RR01}
K.~R. Rajagopal and M.~R{\r u}\v{z}i\v{c}ka.
\newblock Mathematical modeling of electrorheological materials.
\newblock {\em Cont. Mech. and Thermodynamics}, 13(1):59--78, 2001.

\bibitem{Ruz00}
M.~R{\r u}\v{z}i\v{c}ka.
\newblock {\em Electrorheological fluids: modeling and mathematical theory},
  volume 1748 of {\em Lecture Notes in Mathematics}.
\newblock Springer-Verlag, Berlin, 2000.

\bibitem{Zhi97}
V.~V. Zhikov.
\newblock On some variational problems.
\newblock {\em Russian J. Math. Phys.}, 5(1):105--116 (1998), 1997.

\end{thebibliography}

%\bibliographystyle{abbrv}
%\bibliography{references}

\end{document}
